\documentclass[12pt, a4paper]{article}

\usepackage[a4paper, top = 1 in, bottom = 1 in, left = 1 in, right = 1 in]{geometry}

\usepackage{amsmath, amssymb, amsfonts}
\allowdisplaybreaks[4]

\usepackage{enumerate}
\usepackage{multirow}
\usepackage{hhline}
\usepackage{array}
\usepackage{longtable}
\usepackage{graphicx}
\usepackage{tabularray}
\usepackage{undertilde}
\usepackage{xcolor}
\usepackage{dingbat}
\usepackage{fontawesome5}
\usepackage[colorlinks=true, linkcolor=blue, urlcolor=red]{hyperref}
\usepackage{tasks}
\usepackage{bbding}
\usepackage{twemojis}
% how to use bull's eye ----- \scalebox{2.0}{\twemoji{bullseye}}
\usepackage{customdice}
% how to put dice face ------ \dice{2}
\usepackage{fontspec}
\usepackage{fancyhdr}
\usepackage{lastpage}



\pagestyle{fancy}
\fancyhf{}
\fancyfoot[C]{Page \thepage\ of \pageref{LastPage}}
\renewcommand{\headrulewidth}{0pt}



\title{MSMS 401 : Bayesian Statistics and Computation}
\author{{\fontsize{25}{20} \benyorithafont Ananda Biswas}}
\date{}

\newfontface\gloriahallelujahfont{GLORIAHALLELUJAH-REGULAR.TTF}

\newfontface\benyorithafont{Benyoritha-G334D.otf}


\begin{document}

\maketitle

\section*{\faArrowAltCircleRight[regular] \textcolor{blue}{Problem}}

Consider a random sample of size $n$ from $\text{Exp}(\lambda)$, $\lambda$ is the rate parameter. Take a conjugate prior for $\lambda$. Obtain the posterior distribution of $\lambda$, hence obtain Bayes estimator of $\lambda$ assuming squared-error loss function.

\section*{\faArrowAltCircleRight[regular] \textcolor{blue}{Solution}}

Let $x_1, x_2, \ldots, x_n \overset{iid}{\sim} f(x \mid \lambda)$ where

\begin{align*}
f(x \mid \lambda) = \lambda e^{-\lambda x}.
\end{align*}

Consider \textit{a priori} $\lambda \sim \text{Gamma}\left( \text{shape} = \alpha, \text{rate} = \beta \right)$ \textit{i.e.}

\begin{align*}
\pi(\lambda; \alpha, \beta) = \dfrac{\beta^\alpha}{\Gamma(\alpha)} \lambda^{\alpha-1}e^{-\beta\lambda}, \; \lambda > 0.
\end{align*}

The likelihood for $\lambda$ is

\begin{align*}
f(\mathbf{x} \mid \lambda) &= \prod \limits_{i = 1}^{n} \lambda e^{-\lambda x_i} \\[1em]
&= \lambda^n \cdot \exp \left\{ - \lambda \sum \limits_{i = 1}^{n} x_i \right\}
\end{align*}

We compute the posterior distribution of $\lambda$ as
\begin{align*}
p(\lambda \mid \mathbf{x}) &= \dfrac{f(\mathbf{x} \mid \lambda) \, \cdot \, \pi(\lambda)}{\displaystyle \int\limits_{0}^{\infty} f(\mathbf{x} \mid \lambda) \, \cdot \, \pi(\lambda) \, d\lambda} \\[1em]
&= \dfrac{\lambda^n \cdot \exp \left\{ - \lambda \sum \limits_{i = 1}^{n} x_i \right\} \cdot \dfrac{\beta^\alpha}{\Gamma(\alpha)} \lambda^{\alpha-1}e^{-\beta\lambda}}{\displaystyle{\int \limits_{0}^{\infty} \lambda^n \cdot \exp \left\{ - \lambda \sum \limits_{i = 1}^{n} x_i \right\}} \cdot \dfrac{\beta^\alpha}{\Gamma(\alpha)} \lambda^{\alpha-1}e^{-\beta\lambda} \, d\lambda} \\[1em]
&= \dfrac{\lambda^{n + \alpha - 1} \cdot \exp \left\{ - \lambda \left(\sum \limits_{i = 1}^{n} x_i + \beta \right) \right\}}{\displaystyle{\int \limits_{0}^{\infty} \lambda^{n + \alpha - 1} \cdot \exp \left\{ - \lambda \left(\sum \limits_{i = 1}^{n} x_i + \beta \right) \right\} \, d\lambda}} \\[1em]
&= \dfrac{\left( \sum \limits_{i = 1}^{n} x_i + \beta \right)^{n + \alpha}}{\Gamma(n + \alpha)} \cdot \lambda^{n-1} \cdot \exp \left\{ - \lambda \sum \limits_{i = 1}^{n} x_i \right\}.
\end{align*}

which we identify as the PDF of Gamma$\left(\text{shape} = \alpha + n, \text{rate} = \beta + \sum \limits_{i = 1}^{n} x_i \right)$.

\vspace{0.5cm}

$\therefore$ \textit{a posteriori} $\lambda \mid \mathbf{x} \sim \text{Gamma}\left(\alpha + n, \beta + \sum \limits_{i = 1}^{n} x_i \right)$.

\vspace{0.5cm}

{\fontsize{12}{20} \gloriahallelujahfont Under squared error loss function, Bayes estimate is the posterior mean.}

\vspace{0.5cm}

Thus,

\begin{align*}
\hat{\lambda}_{Bayes} = \dfrac{\alpha + n}{\beta + \sum \limits_{i = 1}^{n}  x_i}.
\end{align*}

\end{document}